\documentclass{article}



\usepackage{arxiv}

\usepackage[utf8]{inputenc} % allow utf-8 input
\usepackage[T1]{fontenc}    % use 8-bit T1 fonts
\usepackage{hyperref}       % hyperlinks
\usepackage{url}            % simple URL typesetting
\usepackage{booktabs}       % professional-quality tables
\usepackage{amsfonts}       % blackboard math symbols
\usepackage{nicefrac}       % compact symbols for 1/2, etc.
\usepackage{microtype}      % microtypography
\usepackage{lipsum}		% Can be removed after putting your text content
\usepackage{graphicx}
\usepackage{amsmath}
\usepackage{witharrows} % Equation arrows
\usepackage{cleveref} % Better equation management
\usepackage{doi}
\usepackage{float}
\usepackage{subcaption}
\usepackage{dirtytalk}

% Tikz for graphics and graph drawing
\usepackage{tikz}
\usetikzlibrary{graphs, graphdrawing, graphs.standard, quotes, decorations.pathmorphing, arrows.meta, matrix,decorations.pathreplacing}
\usepackage{listofitems} % For the NN figures
\usepackage[outline]{contour} % glow around text
\contourlength{1.4pt}

\usegdlibrary {force}

% Plots
\usepackage{pgfplots}

\usepackage[backend=biber, style=authoryear, citestyle=authoryear, maxcitenames=2, mincitenames=1, uniquename=false]{biblatex}
\addbibresource{references.bib}
\let\cite\autocite
\crefname{figure}{Figure}{Figures}
\crefname{table}{Table}{Tables}

% Advanced glossary and acronym management
\usepackage[acronym, automake]{glossaries-extra}
\setabbreviationstyle[acronym]{long-short}
\loadglsentries{Acronyms}
\makeglossaries


\title{AmesFormer: A Graph Transformer Neural Network for Mutagenicity Prediction}

%\date{September 9, 1985}	% Here you can change the date presented in the paper title
%\date{} 					% Or removing it

\author{ \href{https://orcid.org/0000-0000-0000-0000}{\includegraphics[scale=0.06]{orcid.pdf}\hspace{1mm}Luke A.~Thompson}\thanks{Use footnote for providing further
		information about author (webpage, alternative
		address)---\emph{not} for acknowledging funding agencies.} \\
	School of Pharmacy\\
	The University of Sydney\\
	Sydney, Australia \\
	\texttt{Placeholder@Nothing.com} \\
	%% examples of more authors
	\And
	\href{https://orcid.org/0000-0000-0000-0000}{\includegraphics[scale=0.06]{orcid.pdf}\hspace{1mm}Josiah ?.~Evans} \\
	Independent Researcher\\
	% Mount-Sheikh University\\
	Seoul, Republic of Korea \\
	\texttt{Placeholder@Nothing.com} \\
    \And
	\href{https://orcid.org/0000-0000-0000-0000}{\includegraphics[scale=0.06]{orcid.pdf}\hspace{1mm}Slade T.~Matthews} \\
	School of Pharmacy\\
	The University of Sydney\\
	Sydney, Australia \\
	\texttt{Placeholder@Nothing.com} \\
	%% \AND
	%% Coauthor \\
	%% Affiliation \\
	%% Address \\
	%% \texttt{email} \\
	%% \And
	%% Coauthor \\
	%% Affiliation \\
	%% Address \\
	%% \texttt{email} \\
	%% \And
	%% Coauthor \\
	%% Affiliation \\
	%% Address \\
	%% \texttt{email} \\
}

% Uncomment to remove the date
%\date{}

% Uncomment to override  the `A preprint' in the header
%\renewcommand{\headeright}{Technical Report}
%\renewcommand{\undertitle}{Technical Report}
\renewcommand{\shorttitle}{AmesFormer}

%%% Add PDF metadata to help others organize their library
%%% Once the PDF is generated, you can check the metadata with
%%% $ pdfinfo template.pdf
\hypersetup{
pdftitle={AmesFormer},
pdfsubject={q-bio.NC, q-bio.QM},
pdfauthor={Luke A.~Thompson, Josiah ?.~Evans},
pdfkeywords={Ames, Mutagenicity, GNN, Transformer},
}

\begin{document}
\maketitle

\begin{abstract}
	The Ames mutagenicity test is a gold standard assay for the safety assessment of new chemicals. However, many \textit{in silico} models rely on challenging-to-interpret ensemble strategies and molecular fingerprint data which neglects gestalt molecular structure. To improve upon these models, we propose AmesFormer, a graph transformer neural network which shows state-of-the-art single model performance for the prediction of Ames mutagenicity. We briefly review the current state of Ames modelling with a focus on graph neural networks. We then discuss the architecture of AmesFormer with specific reference to the generalisation of transformer positional encodings to chemical data and efforts to make AmesFormer accessible to academic researchers. We benchmark AmesFormer using a standardised dataset against 22 other Ames models, achieving near-best-in-class performance. We then quantify the uncertainty of our performance metrics using a unique Bayesian approach. We support our findings with reference to other models from the literature and with theoretical developments in discrete mathematics and graph theory. Overall, we present a high-performance, accessible, and open-source computational model for Ames mutagenicity, with significant potential for regulatory and drug development applications. 
\end{abstract}


% keywords can be removed
\keywords{Ames \and Mutagenicity \and Neural Networks \and GNNs}


\section{Introduction}
\subsection{The Ames Assay}
The Ames test is a widely-used \textit{in vitro} mutagenicity assay essential to drug development. Ames data is explicitly required by many regulatory guidelines, such as \gls{ich} guideline S2 (R1) \cite{ich_ich_2013}. It thus represents a high bar to market access for pharmaceuticals, with many compounds rejected early in the case of an Ames-positive outcome \cite{honma_improvement_2019}.

In the Ames assay, a histidine-deficient substrate is inoculated with strain of auxotrophic mutant histidine-dependent Salmonella typhimirium \cite{ames_improved_1973}. This dependence is introduced via a mutation at the histidine operon \cite{ames_improved_1973}. A suspected mutagenic compound is then introduced. If the mutagen-containing substrate shows significantly greater colony growth than a control, the test-molecule has reversed the histidine-dependence mutation and is mutagenic. This simplicity has resulted in the Ames test having an excellent inter-laboratory replicability of 85\% \cite{kamber_comparison_2009}.

Different \textit{S. typhimurium} strains detect various mutagenicity mechanisms, which can be categorized into substitution mutations (SNPs) or frameshift mutations \cite{lui_mechanistic_2023}. The Ames test also includes S9 rodent liver homogenate, enabling the detection of mutagens that require metabolic activation \cite{maron_revised_1983, ames_improved_1973}.

\subsection{QSAR Modelling}
As the Ames test costs approximately \$2000 per chemical and the CAS reigstry grows by over 4000 molecules daily, the cost of \textit{in vitro} screening for every new chemical quickly becomes prohibitive \cite{honma_improvement_2019}.

To address this problem, \textit{in silico} \gls{qsar} models have been developed to provide cheaper and higher-throughput methods of Ames screening \cite{furuhama_evaluation_2023}. The use of such models is recommended within \gls{ich} guideline M7 (R1) for the control of mutagenic impurities in pharmaceuticals and the European Union's \gls{reach} agreement \cite{european_communities_regulation_2006, ich_assessment_2017, honma_improvement_2019}. \Gls{aicis} and \gls{fda} have also provided guidelines for implementing \textit{in silico} toxicity evaluation methods \cite{aicis_guide_2022, han_fda_2023}.

\subsection{Current Models}
The input to many Ames \gls{qsar} models are \glspl{mf}. \gls{mf}, as seen in \cref{fig:ndma_fp}, are vectorised binary representations of a molecule’s chemical features \cite{capecchi_one_2020}. 

\begin{figure}[H]
    \centering
    \includegraphics[width=0.35\linewidth]{Images/NDMA Fingerprint.png}
    \caption{An illustrative examples of the hashing of different \gls{ndma} substructrures into the \gls{mf} bit vector.}
    \label{fig:ndma_fp}
\end{figure}

Two fingerprinting methodologies are commonly applied to molecules within the Lipinski Limits. 
Structural key \glspl{mf}, such as \gls{maccs} keys seen in \cref{fig:maccs}, encode the presence of a set of predefined chemical substructures into a binary vector \cite{durant_reoptimization_2002, seo_development_2020}.
Hash \gls{mf} \textit{de novo} numerically encode non-predefined substructures. The archetypal \gls{morgan} \gls{mf} shown in \cref{fig:morgan} hashes the local environment (of arbitrary size) around each atom and encodes this into the bit vector representation \cite{rogers_extended-connectivity_2010}.

\begin{figure}[H]
    \centering
    \begin{subfigure}[t]{0.48\textwidth}
        \centering
        \begin{align*}
            \text{MACCS} &= \begin{bmatrix}
               x_{1} \\
               x_{2} \\
               \vdots \\
               x_{166}
             \end{bmatrix}
            \rightarrow \{0, 1\}^{166}
        \end{align*}
        \caption{\gls{maccs} fingerprint are bit vector encoding the presence (1) or absence (0) of a set of 166 predefined molecular substructures. The index of each substructure is unique.}
        \label{fig:maccs}
    \end{subfigure}
    \hfill
    \begin{subfigure}[t]{0.48\textwidth}
        \centering
        \begin{align*}
            \text{ECFP} &= \begin{bmatrix}
               x_{1} \\
               x_{2} \\
               \vdots \\
               x_{n}
             \end{bmatrix}
            \rightarrow \{0, 1\}^n
        \end{align*}
        \caption{Morgan fingerprints encode an arbitrarily sized set of non-predefined molecular substructures into a bit vector of arbitrary length.}
        \label{fig:morgan}
    \end{subfigure}
    \caption{Different molecular fingerprinting techniques and their dimensionalities.}
\end{figure}

\textbf{Number} of the 22? Ames models presented in the Second Ames International Challenge by \textcite{furuhama_evaluation_2023} utilise \gls{mf}-type input data.



Structural Key


\section{Methods}
\subsection{Datasets}
We train AmesFormer on the dataset

\section{Headings: first level}
\label{sec:headings}

\lipsum[4] See Section \ref{sec:headings}.

\subsection{Headings: second level}
\lipsum[5]
\begin{equation}
	\xi _{ij}(t)=P(x_{t}=i,x_{t+1}=j|y,v,w;\theta)= {\frac {\alpha _{i}(t)a^{w_t}_{ij}\beta _{j}(t+1)b^{v_{t+1}}_{j}(y_{t+1})}{\sum _{i=1}^{N} \sum _{j=1}^{N} \alpha _{i}(t)a^{w_t}_{ij}\beta _{j}(t+1)b^{v_{t+1}}_{j}(y_{t+1})}}
\end{equation}

\subsubsection{Headings: third level}
\lipsum[6]

\paragraph{Paragraph}
\lipsum[7]



\section{Examples of citations, figures, tables, references}
\label{sec:others}

\subsection{Figures}
\lipsum[10]
See Figure \ref{fig:fig1}. Here is how you add footnotes. \footnote{Sample of the first footnote.}
\lipsum[11]

\begin{figure}
	\centering
	\fbox{\rule[-.5cm]{4cm}{4cm} \rule[-.5cm]{4cm}{0cm}}
	\caption{Sample figure caption.}
	\label{fig:fig1}
\end{figure}

\subsection{Tables}
See awesome Table~\ref{tab:table}.

The documentation for \verb+booktabs+ (`Publication quality tables in LaTeX') is available from:
\begin{center}
	\url{https://www.ctan.org/pkg/booktabs}
\end{center}


\begin{table}
	\caption{Sample table title}
	\centering
	\begin{tabular}{lll}
		\toprule
		\multicolumn{2}{c}{Part}                   \\
		\cmidrule(r){1-2}
		Name     & Description     & Size ($\mu$m) \\
		\midrule
		Dendrite & Input terminal  & $\sim$100     \\
		Axon     & Output terminal & $\sim$10      \\
		Soma     & Cell body       & up to $10^6$  \\
		\bottomrule
	\end{tabular}
	\label{tab:table}
\end{table}

\subsection{Lists}
\begin{itemize}
	\item Lorem ipsum dolor sit amet
	\item consectetur adipiscing elit.
	\item Aliquam dignissim blandit est, in dictum tortor gravida eget. In ac rutrum magna.
\end{itemize}

\printbibliography
\printglossaries

\end{document}
